\documentclass[11pt,a4paper]{article}
\usepackage[a4paper,margin=1.9cm]{geometry}
\usepackage[T1]{fontenc}
\usepackage[utf8]{inputenc}
\usepackage{amsmath,amssymb,amsthm,mathtools}
\usepackage{xcolor}
\usepackage[most]{tcolorbox}
\usepackage{enumitem}
\usepackage{booktabs}
\usepackage{array}
\usepackage{graphicx}
\usepackage{caption}
\usepackage{float}
\usepackage[colorlinks=true,linkcolor=teal,urlcolor=blue,citecolor=teal]{hyperref}

\definecolor{cQ}{HTML}{1F4E79}      \definecolor{cQBg}{HTML}{F0F5FA}
\definecolor{cKey}{HTML}{1B5E20}    \definecolor{cKeyBg}{HTML}{E8F5E9}
\definecolor{cTrap}{HTML}{B71C1C}   \definecolor{cTrapBg}{HTML}{FFEBEE}
\definecolor{cWhy}{HTML}{4A148C}    \definecolor{cWhyBg}{HTML}{F3E5F5}
\definecolor{cThm}{HTML}{E65100}    \definecolor{cThmBg}{HTML}{FFF3E0}
\definecolor{cNew}{HTML}{AD1457}    \definecolor{cNewBg}{HTML}{FCE4EC}
\definecolor{cSay}{HTML}{01579B}    \definecolor{cSayBg}{HTML}{E1F5FE}
\definecolor{cNote}{HTML}{616161}

\newtcolorbox{keybox}[1][]{breakable,enhanced,colback=cKeyBg,colframe=cKey,
  boxrule=1pt,left=5pt,right=5pt,top=4pt,bottom=4pt,arc=2pt,
  fonttitle=\bfseries\sffamily,coltitle=black,title={#1}}
\newtcolorbox{whybox}[1][]{breakable,enhanced,colback=cWhyBg,colframe=cWhy,
  boxrule=0.8pt,left=5pt,right=5pt,top=4pt,bottom=4pt,arc=2pt,
  fonttitle=\bfseries\sffamily,title={#1}}
\newtcolorbox{thmbox}[1][]{breakable,enhanced,colback=cThmBg,colframe=cThm,
  boxrule=1pt,left=5pt,right=5pt,top=4pt,bottom=4pt,arc=2pt,
  fonttitle=\bfseries\sffamily,coltitle=black,title={#1}}
\newtcolorbox{trapbox}[1][]{breakable,enhanced,colback=cTrapBg,colframe=cTrap,
  boxrule=0.7pt,left=5pt,right=5pt,top=4pt,bottom=4pt,arc=2pt,
  fonttitle=\bfseries\sffamily,title={#1}}
\newtcolorbox{freshbox}[1][]{breakable,enhanced,colback=cNewBg,colframe=cNew,
  boxrule=1pt,left=5pt,right=5pt,top=4pt,bottom=4pt,arc=2pt,
  fonttitle=\bfseries\sffamily,coltitle=black,title={#1}}
\newtcolorbox{saybox}[1][Say this]{breakable,enhanced,colback=cSayBg,colframe=cSay,
  boxrule=0.8pt,left=5pt,right=5pt,top=4pt,bottom=4pt,arc=2pt,
  fonttitle=\bfseries\sffamily,title={#1}}
\newtcolorbox{qabox}[1][]{breakable,enhanced,colback=cQBg,colframe=cQ,
  boxrule=0.7pt,left=5pt,right=5pt,top=4pt,bottom=4pt,arc=2pt,
  fonttitle=\bfseries\sffamily,title={#1}}

\newcommand{\dd}{\mathrm{d}}
\newcommand{\E}{\mathbb{E}}
\newcommand{\PP}{\mathbb{P}}
\newcommand{\QQ}{\mathbb{Q}}
\newcommand{\F}{\mathcal{F}}
\newcommand{\half}{\tfrac{1}{2}}
\newcommand{\si}{\sigma_{\mathrm{imp}}}
\newcommand{\sr}{\sigma_{\mathrm{real}}}
\captionsetup{font=small,labelfont=bf}

\usepackage{listings}
\newcolumntype{L}[1]{>{\raggedright\arraybackslash}p{#1}}
\renewcommand{\arraystretch}{1.12}
\setlength{\parindent}{0pt}
\setlength{\parskip}{0.4em}

\lstdefinestyle{py}{
  language=Python, basicstyle=\ttfamily\scriptsize,
  keywordstyle=\color{cQ}\bfseries, commentstyle=\color{cNote}\itshape,
  stringstyle=\color{cKey}, numbers=left, numberstyle=\tiny\color{cNote},
  numbersep=7pt, frame=leftline, framerule=1.2pt, rulecolor=\color{cQ},
  xleftmargin=13pt, breaklines=true, showstringspaces=false,
  aboveskip=6pt, belowskip=6pt,
}

\title{\vspace{-1.5cm}\textbf{The Stochastic Calculus project, explained}\\
\large Volatility as a quadratic variation on the NIFTY~50\\[4pt]
\normalsize What it is, where the idea came from, what data we used and how we got it,\\
the theory behind every step, and which parts came out of class}
\author{MAC F244 --- Prof.\ Mayank Goel \textbullet\ BITS Goa \textbullet\ companion to \texttt{SC\_Report.pdf}}
\date{15 November 2026}

\begin{document}
\maketitle
\vspace{-0.6cm}

\begin{keybox}[How to use this document]
\texttt{SC\_Report.pdf} is the four-page thing you submit. \textbf{This} is the thing you read
before you have to defend it. It assumes you have not looked at the code and explains every
decision, every number and every figure, in the order a person would actually ask about them.
\smallskip
\textbf{If you have twenty minutes}, read \S1 (what it is), \S3 (the 1973 question),
\S7 (which methods were used) and \S13 (what every member must be able to say).
\textbf{If you have an hour}, read all of it. \S15 is the viva plan. If you have an hour, read all of it.
\end{keybox}

\tableofcontents

% =====================================================================
\section{Purpose, and the project in one page}

\textbf{The purpose.} Black--Scholes needs five inputs, and four of them are directly
observable. The fifth, $\sigma$, is not --- the derivation done in class simply assumes it. The
purpose of this project is to show that It\^o calculus removes that asymmetry, because the same
identity $(\dd W)^2=\dd t$ that puts $\half\sigma^2S^2V_{SS}$ into the Black--Scholes PDE also
makes $\sigma$ the quadratic variation of $\log S$ --- and quadratic variation is computable
from an observed path. Volatility stops being a free parameter and becomes a measurement.
\smallskip
\textbf{In one sentence.} This project verifies that claim, measures NIFTY volatility that way,
shows the market's own volatility quote is systematically too high, and then uses It\^o's lemma
to establish exactly what that error is worth in rupees.

\textbf{What the project does.} Three stages, each a prerequisite for the next.
\begin{itemize}[leftmargin=*,topsep=1pt,itemsep=2pt]
\item \textbf{Verification --- establishing that the instrument is accurate.} We derive that
      $RV_n=\sum(\Delta\log S)^2$ should converge to $\int_0^T\sigma_t^2\dd t$, and that its
      error should have standard deviation $\sqrt{(2T/n)\int\sigma^4\dd t}$. Then we simulate
      paths where we \emph{chose} $\sigma(t)$, so we know the right answer, and check.
      \textbf{Result:} slope $-0.4986$ vs theoretical $-0.5$, and the dispersion matches the
      formula to within $2.1\%$ across a $128\times$ range of sampling frequency.
\item \textbf{Measurement --- characterising NIFTY volatility.} The verified instrument is
      applied to 11 years of NIFTY~50. \textbf{Result:} volatility ranges $5.4\%$--$88.5\%$, is massively
      autocorrelated (Ljung--Box $Q(10)=1752$ against a critical value of $23.2$), and rises
      when the market falls (correlation $-0.200$). Constant-$\sigma$ GBM is wrong, and we can
      say precisely how.
\item \textbf{Valuation --- pricing the market's forecast error.} Realised volatility is
      regressed on India VIX. \textbf{Result:} slope $0.664$ not $1$
      ($t_{b=1}=-3.70$) --- the market \emph{overreacts}. The resulting gap is $+2.50$
      volatility points. It\^o's lemma says a delta-hedged short option converts that gap into
      P\&L, and on 138 months the It\^o prediction tracks realised P\&L at
      \textbf{correlation $+0.949$}.
\end{itemize}

\begin{whybox}[Why the three stages have to come in this order]
They are not three topics that happen to share a dataset --- each one earns the right to the
next. \textbf{Measurement is worthless without verification}: if you had not established that
$RV_n$ estimates $\int\sigma^2\dd t$ accurately, then ``NIFTY volatility ranges from $5\%$ to
$88\%$'' could just be estimator noise rather than a fact about the index.
\textbf{Valuation is worthless without measurement}: comparing India VIX against realised
volatility requires that the realised side actually means something.
\smallskip
And the payoff runs back the other way. The P\&L identity holds on real NIFTY paths precisely
\emph{because} the only thing it needs from the model is the quadratic variation --- the one
object the first stage proved we can measure. That is also the answer to the hardest question
you will be asked; see \S\ref{sec:defence}.
\end{whybox}

% =====================================================================
\section{Where the idea came from}

\begin{qabox}[Purpose: why this topic, and what it is meant to establish]
The course ends at Black--Scholes. Sitting inside Black--Scholes is a parameter, $\sigma$,
that the derivation treats as known and the real world does not supply. I wanted to take the
one tool the course actually gives you for that problem --- the fact that
$[\log S]_T=\int\sigma^2\dd t$ --- and push it as far as it goes: first to check it is true,
then to measure something real with it, then to show it has consequences in money. The project
is therefore not an application \emph{of} the course bolted onto finance; the course result
\emph{is} the project.
\end{qabox}

The honest intellectual path was this. Prof.\ Goel's post-midterm arc is
\[
\text{Brownian motion}\ \to\ \text{quadratic variation}\ \to\ \text{It\^o integral}
\ \to\ \text{It\^o's formula}\ \to\ \text{GBM}\ \to\ \text{Black--Scholes}.
\]
The pivot of that arc is the lecture where $(\dd W)^2=\dd t$ appears, and the consequence
drilled in class is $\tfrac12\sigma^2 S^2 V_{SS}$ appearing in the PDE. What is usually
\emph{not} drilled is the mirror-image consequence: the same identity, read backwards, says
that if you have a path you can recover $\int\sigma^2\dd t$ from it by summing squared
increments. That is the entire field of realised-volatility measurement, and it falls directly
out of one line of the course. Starting there guaranteed the project would be about the course
rather than decorated with it.

\begin{keybox}[What this fixed about my previous project]
My Mathematical Modelling project (MAC F314) lost marks because its supporting experiments ran
on synthetic data when the brief asked for cited, real datasets. So the first constraint here
was non-negotiable: \textbf{real, cited, reproducible market data}, with the simulation used
\emph{only} where simulation is genuinely the right tool --- the verification stage, where you need to know the true
answer in order to check an estimator against it.
\end{keybox}

% =====================================================================
\section{``Why are you using a model from 1973?''}
\label{sec:1973}

This is the single most likely hostile question, and it has a complete answer. There are four
parts to it; any one of them would do, and together they close the question.

\begin{keybox}[Part 1 --- The project does not claim Black--Scholes is right. It measures how
wrong it is.]
Nothing here argues that NIFTY follows geometric Brownian motion. The project
\textbf{demonstrates the opposite}, quantitatively, and that demonstration is a result rather
than an embarrassment:
\begin{itemize}[leftmargin=*,topsep=1pt,itemsep=0pt]
  \item excess kurtosis $19.3$ against $0$ for a Gaussian;
  \item volatility ranging $5.4\%$--$88.5\%$ when the model says it is a constant;
  \item Ljung--Box $Q(10)=1752$ on squared returns against a critical value of $23.2$;
  \item a leverage correlation of $-0.200$ in a model that is symmetric in $W$;
  \item and the market's own Black--Scholes-implied volatility failing an unbiasedness test
        with slope $0.664$ instead of $1$.
\end{itemize}
Black--Scholes is the \textbf{null hypothesis}, not the conclusion. The variance risk premium
is \emph{defined} as a deviation from it --- you need the 1973 baseline in order to say what
the anomaly is a deviation \emph{from}.
\end{keybox}

\begin{keybox}[Part 2 --- The framework is not dated; it is the language all the modern models
are written in]
It\^o calculus is not one model among several. Heston (1993), SABR (2002), rough volatility
(2014), Bates jump-diffusion, local volatility --- every one of them is an SDE driven by
Brownian motion, and every one is analysed with It\^o's lemma and quadratic variation. Calling
It\^o dated is like calling differentiation dated because Newton is from 1687.
\smallskip
More to the point: \textbf{the central object of this project, $[\log S]_T=\int\sigma_t^2\dd t$,
is model-free}. Quadratic variation is a property of the \emph{path}. It does not care whether
the data-generating process is GBM, Heston, or something with jumps --- which is exactly why
the P\&L identity still holds at $r=0.949$ on data that violates GBM badly.
\end{keybox}

\begin{keybox}[Part 3 --- The measurement machinery is from the 2000s, not the 1970s]
\begin{table}[H]
\centering\small
\begin{tabular}{L{6.2cm}L{1.9cm}L{6.9cm}}
\toprule
\textbf{Tool actually used here} & \textbf{Date} & \textbf{What it is} \\
\midrule
Realised variance as an estimator of integrated variance & 2001--02 &
Andersen--Bollerslev--Diebold--Labys; Barndorff-Nielsen--Shephard. The foundation of modern
volatility econometrics. \\
Model-free implied volatility (the VIX construction) & 2000--03 &
Britten-Jones--Neuberger; CBOE's 2003 methodology revision. \textbf{India VIX itself launched
in 2008.} \\
Bipower variation \& the jump test & 2004--05 &
Barndorff-Nielsen--Shephard; Huang--Tauchen. Used in Section~\ref{sec:robustexp}. \\
Variance risk premium as an asset-pricing object & 2009 &
Carr--Wu; Bollerslev--Tauchen--Zhou. \\
\bottomrule
\end{tabular}
\end{table}
\vspace{-6pt}
Only two ingredients are actually from the 1970s: the Black--Scholes pricing formula, used as
a \emph{calculator} for $\Gamma$ and $\Delta$, and the Boyle--Emanuel (1980) convergence
result. Everything that does the analytical work postdates 2000.
\end{keybox}

\begin{saybox}[Part 4 --- the one-sentence answer, if you only get one sentence]
``I am not using Black--Scholes as a theory of prices --- I am using It\^o's identity
$[\log S]=\int\sigma^2\dd t$, which is model-free, to \emph{measure} volatility; Black--Scholes
is only the null hypothesis I measure the market's deviation from, and the paper's result is
precisely that it is rejected, with slope $0.664$ instead of $1$.''
\end{saybox}

\begin{trapbox}[The follow-up that usually comes next]
\textbf{``Then why not use Heston or a rough volatility model?''}
\smallskip
Because they would not change the answer and would weaken the argument. The P\&L
identity~\eqref{eq:pnl2} holds for \emph{any} continuous semimartingale --- Heston included ---
so fitting Heston would add parameters, add estimation error, and leave the identity exactly as
it is. The project's strength is that its central result needs \textbf{no} assumption about the
volatility process beyond continuity, and the jump test in Section~\ref{sec:robustexp} is what
checks that one remaining assumption. Adding a stochastic-volatility model would be
\emph{assuming} what we instead \emph{measured}.
\end{trapbox}

\section{The theory --- and exactly how much of it is class material}
\label{sec:theory}

This is the section to read before a viva. The project is deliberately built so that its spine
is course material and its extensions are clearly labelled.

\subsection{What the course gave us}

\begin{thmbox}[Covered in class, before the midterm]
Finite probability spaces; binomial pricing (European, American, Asian, lookback, knock-out);
$\sigma$-algebras and filtrations; conditional expectation and its five properties;
martingales; the First Fundamental Theorem of asset pricing; Lebesgue integration; Jensen's
inequality; \textbf{Radon--Nikod\'ym and change of measure}; modes of convergence; MCT and DCT;
and --- directly used here --- \textbf{total and quadratic variation}.
\end{thmbox}

\begin{thmbox}[Covered in class, after the midterm]
Brownian motion as a continuous-time object; \textbf{quadratic variation $[W]_T=T$ almost
surely}; the \textbf{It\^o integral} and the It\^o isometry; \textbf{It\^o's (It\^o--Doeblin)
formula}; \textbf{geometric Brownian motion}; \textbf{Black--Scholes--Merton}, including the
delta-hedging derivation of the PDE.
\end{thmbox}

\begin{trapbox}[Flagged in the course as forward material --- so, not assumed here]
Girsanov's theorem, the Feynman--Kac formula, the Ornstein--Uhlenbeck process and
multidimensional It\^o. The project is arranged so that \textbf{none of its results depend on
these}. Where the risk-neutral measure appears, it appears only as the observation that $\mu$
has dropped out of the hedging argument by itself --- which the Black--Scholes derivation done
in class already shows, without needing Girsanov to justify it.
\end{trapbox}

\subsection{Ingredient-by-ingredient provenance}

\begin{table}[H]
\centering\small
\begin{tabular}{L{6.4cm}L{3.1cm}L{5.3cm}}
\toprule
\textbf{Ingredient} & \textbf{Where from} & \textbf{Used for} \\
\midrule
$[W]_T=T$, $(\dd W)^2=\dd t$ & class (post-mid) & the definition of volatility, \eqref{eq:qv2} \\
It\^o's formula & class (post-mid) & $\dd\log S$; the P\&L identity \\
It\^o isometry & class (post-mid) & $[\int\sigma\dd W]=\int\sigma^2\dd t$ \\
Geometric Brownian motion & class (post-mid) & the simulation in verification, the model in valuation \\
Black--Scholes PDE \& delta hedging & class (post-mid) & the P\&L identity \eqref{eq:pnl2} \\
Quadratic variation as a concept & class (pre-mid) & the whole framing \\
Conditional expectation, martingales & class (pre-mid) & why the hedge error has zero mean \\
\midrule
$\operatorname{Var}((\Delta W)^2)=2\Delta^2$ and the CLT constant for $RV_n$
  & \textbf{beyond the course} & the sharp check in verification, \eqref{eq:rvse2} \\
Boyle--Emanuel $n^{-1/2}$ discrete-hedging law & \textbf{beyond the course} & the
  discretisation cost \\
Mincer--Zarnowitz forecast regression & \textbf{beyond} (econometrics) & the unbiasedness test \\
Variance risk premium & \textbf{beyond} (asset pricing) & interpreting the valuation \\
Ljung--Box $Q$ statistic & \textbf{beyond} (time series) & measuring clustering \\
\bottomrule
\end{tabular}
\caption{Provenance of every theoretical ingredient. Roughly speaking: the \emph{spine} is
class material and the \emph{instruments} used to test it are not.}
\end{table}

\subsection{Derivation 1: volatility is a quadratic variation}

Start from the model the course gives, but let $\sigma$ depend on time:
\begin{equation}
\dd S_t=\mu_t S_t\,\dd t+\sigma_t S_t\,\dd W_t .
\label{eq:sde2}
\end{equation}
Apply It\^o's formula to $f(S)=\log S$, for which $f'(S)=1/S$ and $f''(S)=-1/S^2$:
\[
\dd(\log S_t)=\frac{1}{S_t}\dd S_t-\frac{1}{2S_t^2}(\dd S_t)^2 .
\]
Now $(\dd S)^2=\sigma_t^2S_t^2(\dd W)^2=\sigma_t^2S_t^2\dd t$, \emph{using the class fact
$(\dd W)^2=\dd t$ and discarding $\dd t\,\dd W$ and $(\dd t)^2$}. So
\[
\dd(\log S_t)=\Big(\mu_t-\tfrac{\sigma_t^2}{2}\Big)\dd t+\sigma_t\,\dd W_t .
\]
Take quadratic variation of both sides. A $\dd t$ term contributes nothing (it is of bounded
variation), and the It\^o isometry gives the quadratic variation of the stochastic integral:
\begin{equation}
[\log S]_T=\Big[\int_0^T\sigma_t\,\dd W_t\Big]_T=\int_0^T\sigma_t^2\,\dd t\ \eqqcolon\ IV .
\label{eq:qv2}
\end{equation}

\begin{whybox}[Read that backwards --- this is the whole project]
Forwards, \eqref{eq:qv2} is a fact about a model. Backwards it is a \textbf{measurement
procedure}: if you observe the path, you can compute its quadratic variation by summing squared
log-increments, and that sum estimates $\int\sigma^2\dd t$. Volatility stops being a free
parameter and becomes an observable. Everything in \S\ref{sec:method} is built on this.
\end{whybox}

\subsection{Derivation 2: how accurate is the measurement?}

The sample version of \eqref{eq:qv2} on a partition into $n$ pieces of length $\Delta=T/n$ is
$RV_n=\sum_{i}(\log S_{t_i}-\log S_{t_{i-1}})^2$. Each increment is
$(\mu-\tfrac12\sigma^2)\Delta+\sigma_i\Delta W_i$, and squaring it, the $\Delta^2$ and
$\Delta\cdot\Delta W$ cross terms are of smaller order, so the leading behaviour is
$\sigma_i^2(\Delta W_i)^2$. Since $\Delta W_i\sim\mathcal{N}(0,\Delta)$, we have
$(\Delta W_i)^2/\Delta\sim\chi^2_1$, whose variance is $2$; hence
$\operatorname{Var}\big((\Delta W_i)^2\big)=2\Delta^2$. Increments are independent, so variances
add:
\[
\operatorname{Var}(RV_n)\approx\sum_i\sigma_i^4\cdot2\Delta^2
=2\Delta\sum_i\sigma_i^4\Delta\ \xrightarrow[n\to\infty]{}\ 2\Delta\int_0^T\sigma_t^4\,\dd t,
\]
\begin{equation}
\boxed{\ \operatorname{sd}(RV_n-IV)\ \simeq\ \sqrt{\frac{2T}{n}\int_0^T\sigma_t^4\,\dd t}\ }
\label{eq:rvse2}
\end{equation}

\begin{keybox}[Why this matters more than it looks]
\eqref{eq:rvse2} predicts a \textbf{number}, not just a rate. Lots of broken estimators would
still show a $-1/2$ slope on a log--log plot, because $-1/2$ is just what independent averaging
gives you. Far fewer would land on the right \emph{constant}. Checking the constant is
therefore a much sharper test, and it is the one the verification stage runs.
\end{keybox}

\subsection{Derivation 3: what a mispriced volatility is worth}

This is the result that connects the measured gap to money, and it is pure class material --- It\^o's formula
plus the Black--Scholes PDE, nothing else.

Sell one call, priced and hedged at the \emph{implied} volatility $\si$, while the market
actually realises $\sigma_t$. Let $V(S,t)$ be the Black--Scholes value at vol $\si$, so by the
PDE done in class
\begin{equation}
V_t+\half\si^2S^2V_{SS}+rSV_S-rV=0 .
\label{eq:bspde}
\end{equation}
Let $X_t$ be the hedger's wealth, holding $\Delta_t=V_S$ shares with the remainder in the bank.
Self-financing means
\[
\dd X_t=\Delta_t\,\dd S_t+r\big(X_t-\Delta_tS_t\big)\dd t .
\]
Track the hedging \emph{shortfall} $Y_t=X_t-V(S_t,t)$. By It\^o,
$\dd V=V_t\dd t+V_S\dd S+\half\sigma_t^2S^2V_{SS}\dd t$, so
\[
\dd Y=\underbrace{\Delta\,\dd S-V_S\,\dd S}_{=\,0\ \text{because }\Delta=V_S}
+\;r(X-V_SS)\dd t-V_t\,\dd t-\half\sigma_t^2S^2V_{SS}\,\dd t .
\]
\textbf{The $\dd W$ terms are gone} --- that is what delta hedging is for. Substituting $V_t$
from \eqref{eq:bspde} and cancelling the $rSV_S$ terms,
\[
\dd Y_t=rY_t\,\dd t+\half S_t^2V_{SS}\big(\si^2-\sigma_t^2\big)\dd t .
\]
Solving this linear ODE in $Y$ with $Y_0=0$ (the option was sold for exactly $V$) gives, with
$\Gamma_t=V_{SS}$,
\begin{equation}
\boxed{\ \Pi_T=\half\int_0^T e^{r(T-t)}S_t^2\,\Gamma_t\,\big(\si^2-\sigma_t^2\big)\dd t\ }
\label{eq:pnl2}
\end{equation}

\begin{saybox}[If an examiner asks you to state what \eqref{eq:pnl2} means]
``Selling an option and delta-hedging it is a bet on variance and on nothing else. The $\dd W$
term cancels by construction, so direction does not matter. What is left is the gap between the
variance you charged for and the variance that happened, accumulated through time and weighted
by dollar gamma $S^2\Gamma$ --- which is just a measure of how much the hedge is exposed at
that moment. If implied exceeds realised, you make money, path by path.''
\end{saybox}

\begin{trapbox}[The point you must not get wrong]
\eqref{eq:pnl2} is \textbf{pathwise}. It is not an expectation, it is not an
asymptotic approximation, and it does \textbf{not} require returns to be normal, independent
or identically distributed. It needs only that $S$ is a continuous semimartingale whose
quadratic variation is $\int\sigma_t^2\dd t$. This is the single most important defensive point
in the project --- see \S\ref{sec:defence}.
\end{trapbox}

% =====================================================================
\section{The data: what it is, where it came from, how we got it}
\label{sec:data}

\subsection{What we used}

\begin{table}[H]
\centering\small
\begin{tabular}{L{2.9cm}L{3.1cm}L{8.6cm}}
\toprule
\textbf{Series} & \textbf{Ticker} & \textbf{What it is} \\
\midrule
NIFTY~50 index & \texttt{\^{}NSEI} & India's benchmark large-cap equity index. Daily
open/high/low/close, auto-adjusted. This is our $S_t$. \\
India VIX & \texttt{\^{}INDIAVIX} & The NSE's own 30-day implied volatility index, computed
from out-of-the-money NIFTY option prices by the variance-swap construction (the same method
as the CBOE VIX). Quoted in annualised \%. This is our $\si$. \\
\bottomrule
\end{tabular}
\end{table}

\textbf{Coverage:} 2{,}903 trading days, 2~January~2015 to 12~November~2026, after intersecting
the two series on common dates. Zero missing values.

\begin{whybox}[Why India VIX is the right instrument, and not a convenience]
We need the market's \emph{quoted} volatility, $\si$. We could have backed it out of a single
option price, but then the answer would depend on which strike we chose --- the skew would
contaminate it. India VIX is \textbf{model-free}: it is built from a strip of OTM options
across strikes to replicate the price of a variance swap, so it estimates
$\E^\QQ[\int\sigma^2\dd t]$ directly rather than inverting Black--Scholes at one point. That is
exactly the object \eqref{eq:pnl2} compares against, so the match is conceptual, not just
convenient.
\end{whybox}

\subsection{How it was extracted}

Yahoo Finance, through the \texttt{yfinance} Python client, version 1.8.0.

\begin{lstlisting}[style=py]
def _download(ticker: str) -> pd.DataFrame:
    import yfinance as yf
    df = yf.download(ticker, start=START, end=END, progress=False,
                     auto_adjust=True, threads=False)
    if df is None or len(df) == 0:
        raise RuntimeError(f"no rows returned for {ticker}")
    if isinstance(df.columns, pd.MultiIndex):   # yfinance >=0.2 shape
        df.columns = df.columns.get_level_values(0)
    return df
\end{lstlisting}

The first call downloads and writes \texttt{data/nsei.csv} and \texttt{data/indiavix.csv};
every call after that reads the CSV.

\begin{keybox}[Why we cache, and why it is not laziness]
Vendors silently revise history. If the report quotes ``excess kurtosis $19.3$'' and someone
re-runs the code six months later against a revised series, they get a different number and
conclude the report is wrong. \textbf{Caching makes the dataset part of the submission.} The
numbers in \texttt{SC\_Report.pdf} are reproducible from \texttt{data/*.csv} forever, and
\texttt{refresh=True} is there if you genuinely want fresh data.
\end{keybox}

\begin{trapbox}[What we tried first, and why we abandoned it]
The first attempt used \textbf{stooq.com}, which historically served plain CSV. It now
responds with a JavaScript proof-of-work bot check instead of data. Defeating a bot check to
scrape a site is not something we were willing to do, so we switched to Yahoo Finance, which
publishes a supported client library for exactly this purpose. Worth mentioning if asked ---
it shows the data source was chosen, not stumbled into.
\end{trapbox}

\subsection{How we know the data is not garbage}

Before any series is returned, it is checked. If a check fails the program stops; it never
reports a number computed from bad data.

\begin{lstlisting}[style=py]
def _validate(frame: pd.DataFrame) -> None:
    assert len(frame) > 2000, f"only {len(frame)} rows; expected a decade"
    assert frame.index.is_monotonic_increasing, "dates out of order"
    assert not frame.index.duplicated().any(), "duplicate dates"
    assert (frame[["open","high","low","spot"]] > 0).all().all(), "non-positive price"
    assert (frame["high"] >= frame["low"]).all(), "high < low"
    assert frame["iv"].between(0.02, 1.50).all(), "implied vol outside a sane band"
    gap = frame.index.to_series().diff().dt.days.dropna()
    assert gap.max() <= 15, f"a {gap.max()}-day hole in the series"
\end{lstlisting}

Each line catches a specific real failure mode: a truncated download, a vendor returning dates
out of order, a double-counted session, a split-adjustment bug producing zero or negative
prices, corrupted OHLC, a decimal-point error in the VIX, and a silent multi-week hole from a
delisting or feed outage.

\subsection{What the data says before we do anything to it}

\begin{table}[H]
\centering\small
\begin{tabular}{lr@{\qquad}lr}
\toprule
Trading days & $2{,}903$ & Annualised realised vol (full sample) & $16.33\%$ \\
First / last & 2015-01-02 / 2026-11-12 & Mean India VIX & $16.54\%$ \\
Missing values & $0$ & \textbf{Excess kurtosis of daily log returns} & $\mathbf{19.28}$ \\
Interest rate used & $6.5\%$ flat & \textbf{Skew of daily log returns} & $\mathbf{-1.27}$ \\
\bottomrule
\end{tabular}
\end{table}

\begin{trapbox}[Do not hide the kurtosis --- lead with it]
Excess kurtosis $19.28$ means the daily returns are \emph{nothing like} normal; a Gaussian has
excess kurtosis $0$. GBM assumes normal log returns. So the model we are using is, as a
description of the data, clearly false. \textbf{Say this before the examiner does}, then give
the answer in \S\ref{sec:defence}. A project that reports this number and explains why its
conclusions survive is far stronger than one that quietly omits it.
\end{trapbox}

% =====================================================================
\section{The method, step by step}
\label{sec:method}

\subsection{Verification: testing the instrument where the answer is known}

You cannot test $RV_n\to IV$ on market data, because on market data nobody knows $IV$. So we
build a world where we do: simulate \eqref{eq:sde2} with a volatility we chose ourselves,
\[
\sigma(t)=0.15+0.10\sin(2\pi t),\qquad T=1,
\]
deliberately \emph{time-varying}, so the test is not accidentally passed by a constant-$\sigma$
special case. Numerical quadrature gives the two constants the theory needs:
\[
\int_0^1\sigma^2\dd t=0.027500,\qquad \int_0^1\sigma^4\dd t=0.001219 .
\]
Then for each $n\in\{21,42,\dots,2688\}$ we simulate 4{,}000 paths, compute $RV_n$ on each,
and compare the distribution of $RV_n-IV$ against \eqref{eq:rvse2}.

\begin{lstlisting}[style=py]
dt   = T / n
grid = np.arange(n) * dt
sig  = sigma_path(grid + 0.5 * dt)          # midpoint rule: 2nd-order accurate
z    = rng.standard_normal((n_paths, n))
dlog = (mu - 0.5 * sig**2) * dt + sig * np.sqrt(dt) * z
rv   = (dlog**2).sum(axis=1)                # this is RV_n, one per path
\end{lstlisting}

\textbf{Result.} Bias never exceeds $1.1\%$ of $IV$ and is under $0.35\%$ for $n\ge42$. The
log--log slope is $-0.4986$ against the predicted $-0.5$. The ratio of measured to theoretical
standard deviation stays in $[0.991,1.021]$ across the whole range --- i.e.\ the formula is
right to about $2\%$ over a $128$-fold change in sampling frequency.

\subsection{Measurement: applying it to the NIFTY~50}

Apply $RV_n$ over rolling 21-day windows (one listed month) and annualise. Then ask what the
resulting series looks like, using three standard diagnostics.

\begin{itemize}[leftmargin=*,topsep=1pt,itemsep=2pt]
\item \textbf{Clustering} --- autocorrelation of $|r_t|$ and $r_t^2$, plus Ljung--Box
      $Q(10)=n(n+2)\sum_{k=1}^{10}\hat\rho_k^2/(n-k)$, which is $\chi^2_{10}$ under the null of
      no structure.
\item \textbf{Leverage} --- $\operatorname{corr}(r_t,\Delta\hat\sigma_t)$, the correlation
      between today's return and the change in measured volatility around it.
\item \textbf{Vol-of-vol} --- standard deviation of daily changes in $\log\hat\sigma$.
\end{itemize}

\textbf{Results, and what each one kills:}
\begin{table}[H]
\centering\small
\begin{tabular}{L{4.6cm}L{3.0cm}L{7.3cm}}
\toprule
\textbf{Measurement} & \textbf{Value} & \textbf{What it rules out} \\
\midrule
Realised vol, mean / sd & $14.17\%$ / $8.08\%$ & --- \\
Realised vol, range & $5.37\%$ -- $88.45\%$ & that any single $\sigma$ describes the sample \\
ACF of $|r_t|$, lag 1 / lag 42 & $+0.256$ / $+0.080$ & fast (exponential) mean reversion in vol \\
Ljung--Box $Q(10)$, squared returns & $\mathbf{1752.1}$ (crit.\ $23.2$) & i.i.d.\ returns, overwhelmingly \\
Ljung--Box $Q(10)$, returns & $59.3$ (crit.\ $23.2$) & \emph{also} rejected --- see the caveat below \\
Leverage correlation & $-0.200$ & symmetry in $W$; explains the option skew \\
Vol-of-vol (annualised) & $99.2\%$ & constant $\sigma$ --- but see the caveat \\
\bottomrule
\end{tabular}
\end{table}

\begin{trapbox}[Two caveats we put in the report rather than hiding]
\textbf{(1)} $Q(10)=59.3$ on the raw returns also exceeds the critical value, driven by lags 5
and 10. So it is \emph{not} true that ``returns are unpredictable and only volatility is''. The
correct statement is that return autocorrelation is small but statistically detectable, an
order of magnitude weaker than the variance structure. Claiming zero would have been wrong, and
an examiner reading the ACF table would catch it.
\smallskip
\textbf{(2)} The $99.2\%$ vol-of-vol is an \textbf{upper bound}, not an estimate. We obtained it
by differencing a 21-day rolling window \emph{daily}; consecutive windows share 20 of 21
observations, so the difference is dominated by measurement noise rather than by genuine
movement in the latent volatility. It is quoted for order of magnitude and nothing in the
project depends on it.
\end{trapbox}

\subsection{Valuation: the market's forecast error}

Take 138 \textbf{non-overlapping} 21-day cycles. For each, record the India VIX on day one and
the realised volatility over the following 21 days. Then run the Mincer--Zarnowitz regression
\[
\sr=a+b\,\si+\varepsilon .
\]
If implied volatility were an unbiased forecast, we would find $a=0$ and $b=1$.

\begin{table}[H]
\centering\small
\begin{tabular}{lrrr}
\toprule
 & estimate & std.\ error & $t$ \\
\midrule
intercept $a$ & $+0.0310$ & $0.0163$ & $+1.91$ \\
slope $b$     & $+0.6637$ & $0.0909$ & $+7.30$ \\
\midrule
\multicolumn{4}{l}{$R^2=0.282$, $n=138$. Test of $b=1$: $t=-3.70$ \ (\textbf{rejected}).} \\
\bottomrule
\end{tabular}
\end{table}

\textbf{Reading it.} India VIX is genuinely informative --- it explains $28\%$ of the variation
in what follows, and the slope is $7.3$ standard errors above zero. But it is \textbf{biased}:
$b=0.66$ means a one-point rise in India VIX is followed, on average, by only two-thirds of a
point of extra realised volatility. The market overreacts to volatility news.

\begin{whybox}[Why non-overlapping cycles, and why this is a real methodological choice]
Using overlapping windows would give us $\sim$2{,}900 observations instead of $138$, and much
prettier standard errors. They would also be \textbf{wrong}: adjacent overlapping windows share
almost all of their data, so the residuals are heavily autocorrelated and the classical OLS
standard errors understate the truth badly --- you would ``discover'' significance that is not
there. We took the smaller honest sample. If asked why $n$ is only 138, this is the answer.
\end{whybox}

\textbf{The premium.} Mean implied $16.66\%$, mean subsequent realised $14.16\%$: a gap of
$\mathbf{+2.50}$ volatility points, positive in $76.8\%$ of the 138 months, $t=+4.04$. In
variance terms $\E[\si^2-\sr^2]=+0.00531$ --- which is precisely the quantity that
\eqref{eq:pnl2} says a delta-hedged short option collects.

\subsection{Closing the loop: from the gap to the money}

Two final tests confirm that \eqref{eq:pnl2} is not just algebra.

\textbf{(i) The cost of hedging discretely.} \eqref{eq:pnl2} assumes continuous rebalancing,
which is impossible. Boyle--Emanuel (1980) predicts that with $n$ rebalances the residual error
has standard deviation $\propto n^{-1/2}$. On simulated GBM with $\sigma$ known exactly --- so
discretisation is the \emph{only} source of error --- we measure a log--log slope of $-0.470$
across all settings, tightening monotonically to $\mathbf{-0.489}$ for $n\ge21$ where the
asymptotic regime actually applies. The mean error is indistinguishable from zero at every $n$:
\textbf{discrete hedging is unbiased, and it is the dispersion that costs you}.

\textbf{(ii) The identity on real paths.} For each of the 138 cycles, evaluate the integral in
\eqref{eq:pnl2} along the realised NIFTY path, and compare it against the actual P\&L of the
hedged position computed independently by the self-financing simulator. They agree at
\textbf{correlation $+0.949$} (predicted $+39.2$ points per cycle against actual $+36.1$).

\begin{table}[H]
\centering\small
\begin{tabular}{lr@{\qquad}lr}
\toprule
Cycles (non-overlapping) & $138$ & Mean P\&L & $+36.06$ pts \\
As \% of premium received & $+11.31\%$ & Standard deviation & $88.30$ \\
Win rate & $74.6\%$ & Cumulative & $+4{,}977$ pts \\
$t$-statistic & $4.80$ & $p$-value & $4.1\times10^{-6}$ \\
\bottomrule
\end{tabular}
\end{table}

% =====================================================================
\subsection{Robustness: the three attacks, answered by computation}
\label{sec:robustexp}

\textbf{(a) ``Your estimator is measuring jumps, not diffusion volatility.''}
Identity~\eqref{eq:qv2} assumes $S$ is \emph{continuous}. If it jumps,
$RV_n\to\int\sigma^2\dd t+\sum(\text{jump})^2$ and the two parts must be separated. The tool is
\textbf{bipower variation},
\[
  BV=\frac{\pi}{2}\sum_{i=2}^{n}|r_i|\,|r_{i-1}| .
\]
\begin{whybox}[Why multiplying neighbours kills jumps --- the idea in one line]
In $RV$ a jump enters \emph{squared}: a return ten times normal contributes $100\times$. In
$BV$ that same large $|r_i|$ is multiplied by its \textbf{neighbour}, which is a normal-sized
return, so it contributes only $10\times$ and only to one term of the sum. As $n\to\infty$ the
finitely many jump terms are averaged away, while the diffusive part survives --- so $BV$
estimates $\int\sigma^2\dd t$ alone, and $RV-BV$ estimates the jumps. The $\pi/2$ is
$\big(\E|Z|\big)^{-2}$ for $Z\sim\mathcal{N}(0,1)$, which rescales $\E|r_i||r_{i-1}|$ back to a
variance.
\end{whybox}
The Barndorff-Nielsen--Shephard ratio statistic (in the Huang--Tauchen form) standardises
$(RV-BV)/RV$ by its asymptotic standard error, which needs the integrated \emph{quarticity},
estimated jump-robustly by tripower variation:
\[
  Z=\frac{(RV-BV)/RV}
        {\sqrt{\big((\tfrac{\pi}{2})^2+\pi-5\big)\tfrac{1}{n}\max\!\big(1,TQ/BV^2\big)}}
   \ \sim\ \mathcal{N}(0,1)\ \text{ under ``no jumps''.}
\]
\textbf{Result:} over the full sample the jump share is $\mathbf{5.60\%}$ and
$Z=+1.617$ --- below the $5\%$ critical value of $1.96$, so the no-jump null is \textbf{not}
rejected. Year by year it is rejected only in 2015 and 2016 (2 of 12 years). About $94\%$ of
NIFTY quadratic variation is diffusive, which is what licenses the continuity assumption.

\textbf{(b) ``Your profit vanishes once you pay to trade.''}
Every rebalance now charges a proportional fee on the notional traded. Three separate things
come out, and keeping them separate is the point:

\begin{table}[H]
\centering\small
\begin{tabular}{L{3.3cm}L{11.9cm}}
\toprule
\textbf{Quantity} & \textbf{What it does, and what it means} \\
\midrule
The \textbf{mean} & \emph{Not} a function of rebalance frequency: the spread across
frequencies is $22.5$ points against a typical standard error of $11.9$. This is a
\textbf{confirmation}, not a nuisance --- discrete hedging is unbiased, so frequency must not
move the mean. \\
The \textbf{dispersion} & Falls monotonically with frequency: standard deviation $230$ points
with two rebalances, $88$ hedging daily. This is the $n^{-1/2}$ law appearing again, now in a
real backtest rather than a simulation. \\
The \textbf{cost drag} & Deterministic, and it grows \textbf{sublinearly} --- a log--log slope
of $+0.30$ against trade count, not $+1$. The reason is worth understanding: hedging twice as
often does \emph{not} double your trading, because each individual rebalance is smaller. \\
\bottomrule
\end{tabular}
\end{table}

\textbf{Result:} hedging daily, the premium is $+32.4$ points at $1$~bp ($t=4.30$), $+28.8$ at
$2$~bps ($t=3.81$), $+17.8$ at $5$~bps ($t=2.33$), and reaches break-even near $10$~bps. Indian
index-futures round trips sit at the low end of that range, so the premium is not an artefact
of ignoring costs.

\begin{trapbox}[Do not claim an ``optimal rebalance frequency'']
It is tempting to read the table and announce that weekly hedging is optimal because it has the
highest net P\&L. \textbf{That would be reading noise.} The means are statistically
indistinguishable at zero cost; the apparent winner changes with sampling. The defensible
statement is about the \emph{drag} and the \emph{dispersion}, both of which are monotone and
both of which have theory behind them.
\end{trapbox}

\textbf{(c) ``This is one lucky sample.''}
Splitting into disjoint sub-periods: $2015$--$2019$ gives $+24.0$ points per cycle over 59
cycles ($t=3.22$); the COVID window $2020$--$2021$ gives $+56.9$ over 23 cycles ($t=2.91$);
$2022$--$2026$ gives $+40.2$ over 56 cycles ($t=2.74$). Every sub-period is positive and
individually significant, and the premium is \emph{larger} when volatility is higher --- which
is what a risk premium is supposed to do.

\section{Exactly which numerical methods were used --- and which were not}
\label{sec:methods}

Every member of the group should be able to answer this section without hesitation. It is
written as a checklist for that reason.

\subsection{What was used}

\begin{table}[H]
\centering\small
\begin{tabular}{L{3.6cm}L{5.4cm}L{6.0cm}}
\toprule
\textbf{Method} & \textbf{Where exactly} & \textbf{Why that one} \\
\midrule
\textbf{Exact log-solution of the SDE} (\emph{not} Euler on $S$) &
\texttt{sc\_vol.rv\_convergence}, \texttt{sc\_hedge.convergence\_study} &
For $\dd S=\mu S\dd t+\sigma S\dd W$ the log transform integrates exactly:
$\Delta\log S=(\mu-\tfrac{\sigma^2}{2})\Delta+\sigma\sqrt{\Delta}\,Z$. Simulating this has
\textbf{zero discretisation error}, so any error we measure is the thing we are studying, not
the simulator. \\
\addlinespace
\textbf{Monte Carlo} & 4{,}000 paths per setting (convergence), 20{,}000 (hedging) &
Fixed seed \texttt{RNG\_SEED=2026}; the quantity of interest is a distribution, not a point. \\
\addlinespace
\textbf{Closed-form Black--Scholes} & \texttt{sc\_model}: price, $\Delta$, $\Gamma$ &
Available analytically, so no PDE solver is needed; verified against the PDE and finite
differences before use. \\
\addlinespace
\textbf{Central finite differences} & \emph{verification only}, \texttt{check\_greeks} &
To confirm the analytic Greeks ($3\times10^{-8}$ agreement). Not used to produce any result. \\
\addlinespace
\textbf{Trapezoidal quadrature} & $\int\sigma^2\dd t$, $\int\sigma^4\dd t$ over $2\times10^5$
points & The integrand is smooth and deterministic; we need the exact target to test against. \\
\addlinespace
\textbf{Riemann sum along the path} & the It\^o integral in \texttt{ito\_prediction} &
It\^o integrals use the \textbf{left endpoint} --- that is what makes them martingales. Our
loop evaluates $\Gamma$ at $t_i$ and multiplies by the increment over $[t_i,t_{i+1}]$. \\
\addlinespace
\textbf{OLS, implemented by hand} & \texttt{sc\_vol.ols}, Mincer--Zarnowitz &
Written out so the standard errors are ours and can be explained; no regression library. \\
\addlinespace
\textbf{Bipower \& tripower variation} & \texttt{sc\_robust}, the jump test &
Jump-robust estimators of $\int\sigma^2\dd t$ and $\int\sigma^4\dd t$. \\
\bottomrule
\end{tabular}
\end{table}

\subsection{What was deliberately \emph{not} used, and why}

\begin{table}[H]
\centering\small
\begin{tabular}{L{3.6cm}L{11.6cm}}
\toprule
\textbf{Method} & \textbf{Why not} \\
\midrule
\textbf{Euler--Maruyama on $S$} & Unnecessary. The log transform is exact, so Euler would
\emph{introduce} an error we would then have to disentangle from the discrete-hedging error we
are trying to measure. \\
\addlinespace
\textbf{Milstein scheme} & It would change nothing. Milstein adds
$\half b\,b'\big((\Delta W)^2-\Delta\big)$, where $b$ is the diffusion coefficient and
$b'=\partial b/\partial x$. In log coordinates $b=\sigma(t)$ does not depend on the state, so
$b'=0$ and \textbf{Milstein reduces exactly to Euler}. This is worth saying out loud --- it
shows the choice of scheme was reasoned, not defaulted. \\
\addlinespace
\textbf{Binomial / trinomial trees} & A discrete-time pricing device. We need a
\emph{continuous} path to accumulate quadratic variation along; a tree would impose a lattice
that has no quadratic variation in the It\^o sense. \\
\addlinespace
\textbf{Finite-difference PDE solvers} & The Black--Scholes PDE has a closed-form solution for
a European call. A grid solver would add error for no benefit. \\
\addlinespace
\textbf{Neural networks, PINNs, any machine learning} & \textbf{None was used anywhere in this
project.} There is no network, no training loop, no \texttt{torch} import. Every number comes
from a closed-form expression, a Monte Carlo simulation, or an ordinary least-squares fit. If
asked directly: \emph{no neural method of any kind appears in this project.} \\
\addlinespace
\textbf{Girsanov's theorem} & Not needed. The drift $\mu$ cancels out of the hedging argument
by itself --- see the derivation in \S\ref{sec:theory}, where the $\dd S$ terms cancel because
$\Delta=V_S$. Girsanov would be required to \emph{price} from scratch; we only ever
\emph{hedge}. It is also forward material in the course. \\
\addlinespace
\textbf{GARCH / stochastic-volatility fitting} & Deliberately avoided. The P\&L identity holds
for any continuous semimartingale, so fitting a volatility model would add parameters and
estimation error while leaving the result unchanged. We \emph{measure} the quadratic variation
instead of \emph{assuming} a process for it. \\
\bottomrule
\end{tabular}
\end{table}

\begin{saybox}[If someone asks ``which method did you use?'' and you have five seconds]
``Exact log-Euler Monte Carlo for the simulations --- exact because in log coordinates the GBM
increment is closed-form --- closed-form Black--Scholes for the Greeks, a left-endpoint Riemann
sum for the It\^o integral, and hand-written OLS for the regression. No PDE solver, no tree,
no neural network.''
\end{saybox}

\section{Where this is actually used in the real world}
\label{sec:real}

The honest framing is that this project reproduces, on Indian index data, a calculation that
several real businesses run every day.

\begin{table}[H]
\centering\small
\begin{tabular}{L{3.9cm}L{11.3cm}}
\toprule
\textbf{Use} & \textbf{How this project's object appears in it} \\
\midrule
\textbf{Variance swaps} & The payoff is literally
$N\times(\text{realised variance}-\text{strike})$, where realised variance is
$\sum(\Delta\log S)^2$ --- \emph{exactly} the $RV_n$ of Section~\ref{sec:method}, annualised.
These trade OTC in billions of notional daily. The entire contract exists because It\^o makes
quadratic variation a well-defined, replicable quantity. \\
\addlinespace
\textbf{VIX and India VIX futures} & India VIX has been published by the NSE since 2008 and
\textbf{India VIX futures have traded on the NSE since February 2014}. Their pricing is the
market's forecast of exactly the quantity we regress in the valuation stage. \\
\addlinespace
\textbf{Option market-maker P\&L attribution} & Every derivatives desk decomposes each day's
P\&L into delta, \emph{gamma}, vega and theta components. The gamma term is
$\half S^2\Gamma\big(\si^2-\sigma_{\text{realised}}^2\big)\Delta t$ --- line for line our
equation~\eqref{eq:pnl2}. A desk that cannot explain its gamma P\&L cannot explain its day. \\
\addlinespace
\textbf{Short-volatility strategies, and their blow-ups} & Selling options and delta-hedging is
a real, widely-run trade. On \textbf{5 February 2018 (``Volmageddon'')} the VIX more than
doubled in a day and the XIV short-volatility ETN lost about $96\%$ of its value and was
liquidated. That is the tail our sub-period analysis is probing, and it is the reason the
report reports a standard deviation of $88$ points next to a mean of $36$. \\
\addlinespace
\textbf{Risk management and margining} & Realised volatility estimated by quadratic variation
feeds value-at-risk, expected shortfall, and exchange margin systems (SPAN and its successors).
When an exchange raises margins after a volatile week, this estimator is why. \\
\addlinespace
\textbf{Indian market context} & India is now the world's largest derivatives market by number
of contracts traded, overwhelmingly NIFTY and BANKNIFTY options. A persistent variance risk
premium of $+2.5$ volatility points is a statement about who is paying whom in that market:
hedgers and retail option \emph{buyers} pay insurance, and \emph{sellers} are compensated for
carrying crash risk. \\
\bottomrule
\end{tabular}
\end{table}

\begin{whybox}[Why the premium exists at all --- the economics, not the mathematics]
A positive variance risk premium is not a free lunch and not a market inefficiency. Options pay
off precisely when markets crash, which is when investors most need money and when wealth is
scarcest. Assets that pay off in bad states are therefore expensive --- their expected return
is \emph{negative}, by exactly the same logic that makes insurance cost more than its actuarial
value. Selling that insurance earns a premium as compensation for bearing the crash risk.
\smallskip
This is the right answer to ``if this makes money, why doesn't everyone do it?'' --- they
\emph{do}, and the compensation survives because the risk is real. Our own numbers show it:
the strategy loses in $25\%$ of months, and the worst single cycle is $-419$ points against a
mean of $+36$.
\end{whybox}

\section{Every figure, and what to look at}

\begin{table}[H]
\centering\small
\begin{tabular}{L{2.6cm}L{6.2cm}L{6.1cm}}
\toprule
\textbf{Figure} & \textbf{What is plotted} & \textbf{What it is there to show} \\
\midrule
1 \texttt{data\_overview} & NIFTY level on a log axis; below it, India VIX against 21-day
realised vol & That the dataset is real and complete, and that implied sits visibly above
realised most of the time --- the premium, before any statistics. \\
\addlinespace
2 \texttt{rv\_convergence} & Left: sd of $RV_n-IV$ against $n$, log--log, measured points over
the theoretical \emph{line}. Right: bias vs $n$ & That the line is a \textbf{prediction, not a
fit}: the points land on it over a $128\times$ range. Right: bias stays inside $\pm2\%$. \\
\addlinespace
3 \texttt{vol\_clustering} & Left: ACF of $r$, $|r|$, $r^2$ with the $\pm1.96/\sqrt T$ band.
Right: scatter of return against change in vol & Left: grey bars (returns) are near the band;
coloured bars tower over it --- size is predictable, sign is not. Right: the downward slope is
the leverage effect. \\
\addlinespace
4 \texttt{implied\_vs\_realised} & Left: realised vs implied, with the $45^\circ$ unbiasedness
line and the fitted line. Right: the monthly gap & Left: the fitted line is visibly
\emph{flatter} than $45^\circ$ --- that is $b=0.66$. Right: bars are mostly positive. \\
\addlinespace
5 \texttt{hedge} + \texttt{ito\_vs\_actual} & The $n^{-1/2}$ law; cumulative P\&L; and It\^o
prediction against realised P\&L & That the theory's two quantitative predictions --- the
convergence rate and the P\&L identity --- both hold on data. \\
\addlinespace
6 \texttt{robustness} & Realised vs bipower volatility by year; net P\&L against cost; mean
P\&L by sub-period & The three objections, pre-answered: orange and blue bars nearly equal
(jumps are small), the cost curves stay above zero at realistic costs, and all three
sub-period bars are positive with error bars clear of zero. \\
\bottomrule
\end{tabular}
\end{table}

\begin{saybox}[The one figure to put on a slide]
The right-hand panel of Figure 5. A scatter of ``what It\^o's lemma predicted this month'' against
``what actually happened this month'', 138 points, hugging the diagonal, $r=+0.949$. It is the
single clearest statement that the theory is doing real work.
\end{saybox}

% =====================================================================
\section{The code}

Six flat files, no package structure, no framework. Every random draw uses a fixed seed
(\texttt{RNG\_SEED = 2026}), so every number reproduces exactly.

\begin{table}[H]
\centering\small
\begin{tabular}{L{2.6cm}L{2.0cm}L{9.9cm}}
\toprule
\textbf{File} & \textbf{Lines} & \textbf{What it does} \\
\midrule
\texttt{sc\_data.py} & $\sim$85 & Downloads, caches and \textbf{validates} both series.
\texttt{nifty\_and\_vix()} returns the aligned frame; \texttt{\_validate()} is the gatekeeper. \\
\texttt{sc\_model.py} & $\sim$120 & Black--Scholes price, delta and gamma, plus
\texttt{check\_pde()} and \texttt{check\_greeks()} which verify the implementation \emph{before}
it is used. \\
\texttt{sc\_vol.py} & $\sim$200 & Verification (\texttt{rv\_convergence}), measurement
(\texttt{vol\_diagnostics}, \texttt{ljung\_box}) and valuation
(\texttt{implied\_vs\_realised}, \texttt{ols}). \\
\texttt{sc\_hedge.py} & $\sim$175 & \texttt{hedge\_once()} --- the self-financing hedge
simulator, with an optional per-rebalance transaction cost; \texttt{convergence\_study()};
\texttt{backtest()}; \texttt{ito\_prediction()}. \\
\texttt{sc\_robust.py} & $\sim$180 & \texttt{bipower\_variation()},
\texttt{tripower\_quarticity()} and \texttt{jump\_test()}; \texttt{cost\_sweep()};
\texttt{subperiods()}. \\
\texttt{sc\_figures.py} & $\sim$250 & All six figures, read from the CSVs the other modules
write, so plots and report always quote the same numbers. \\
\bottomrule
\end{tabular}
\end{table}

\subsection{The one function worth being able to explain}

\begin{lstlisting}[style=py]
def hedge_once(path, strike, rate, price_vol, hedge_vol, ttm_grid, cost_bps=0.0):
    fee     = cost_bps * 1e-4                  # 0.0 reproduces the frictionless case
    premium = float(call_price(path[0], strike, rate, price_vol, ttm[0]))
    delta   = float(call_delta(path[0], strike, rate, hv[0], ttm[0]))
    cash    = premium - delta * path[0]        # sold the option, bought the hedge
    cash   -= fee * abs(delta) * path[0]       # cost of putting the hedge on

    for i in range(1, len(path)):
        dt    = ttm[i-1] - ttm[i]
        cash *= np.exp(rate * dt)              # the bank account grows
        if ttm[i] > 0:
            new   = float(call_delta(path[i], strike, rate, hv[i], ttm[i]))
            cash -= (new - delta) * path[i]    # rebalance: buy/sell the difference
            cash -= fee * abs(new - delta) * path[i]   # cost of that rebalance
            delta = new

    cash += delta * path[-1]                   # liquidate the share position
    cash -= fee * abs(delta) * path[-1]        # cost of liquidating
    cash -= max(path[-1] - strike, 0.0)        # settle the short call
    return cash                                # == the hedge error
\end{lstlisting}

\begin{keybox}[Why this is written as a simulation and not a formula]
Nothing here assumes the answer. It literally sells the option, borrows, buys shares, pays
interest, rebalances, liquidates and settles. The returned cash is whatever is left. That makes
it an \textbf{independent} check on \eqref{eq:pnl2}: the identity is computed one way (an
integral over the path) and the P\&L another way (an accounting simulation), and the two are
then compared. If we had derived the P\&L from the formula, the correlation of $0.949$ would be
circular and worthless.
\end{keybox}

\subsection{What is verified before any result is computed}

\begin{itemize}[leftmargin=*,topsep=1pt,itemsep=1pt]
\item \textbf{PDE residual:} substituting the implemented price into
      $V_t+\half\sigma^2S^2V_{SS}+rSV_S-rV$ gives $-2.212\times10^{-4}$ against $rV=65.584$ ---
      a relative residual of $3.37\times10^{-6}$.
\item \textbf{Greeks:} analytic delta against a central finite difference, $3.26\times10^{-8}$;
      gamma, relative $1.08\times10^{-7}$.
\item \textbf{Put--call parity:} holds to six decimal places.
\end{itemize}

\begin{whybox}[Why bother]
Because a sign error in gamma would still produce a plausible-looking backtest, and we would
never know. These three checks cost twenty lines and remove an entire category of silent
failure. If asked ``how do you know your Black--Scholes is right?'', this is the answer.
\end{whybox}

% =====================================================================
\section{The defence: why the conclusions survive a false model}
\label{sec:defence}

\begin{qabox}[The question a good examiner will ask]
``Your own data says excess kurtosis is $19.3$ and volatility is neither constant nor symmetric.
Geometric Brownian motion assumes normal returns and constant $\sigma$. Your model is wrong.
Why should I believe a correlation of $0.949$?''
\end{qabox}

\textbf{The answer, in three steps.}

\textbf{One.} Identity \eqref{eq:pnl2} is a \emph{pathwise} consequence of It\^o's lemma.
Re-read the derivation in \S\ref{sec:theory}: at no point does it use normality, independence,
or that $\sigma$ is constant. It uses only that $S$ is a continuous semimartingale and that
$(\dd W)^2=\dd t$, i.e.\ that the quadratic variation exists.

\textbf{Two.} The $\sigma_t^2$ that appears in \eqref{eq:pnl2} is a \textbf{realised quadratic
variation along the path that happened}, not a distributional parameter of an assumed law. Fat
tails and clustering change \emph{what number} $\int\sigma_t^2\dd t$ turns out to be on a given
path. They do not break the identity that relates that number to P\&L.

\textbf{Three.} And the verification stage is what licenses step two: it is precisely the
demonstration that
$\int\sigma_t^2\dd t$ can be measured accurately from a path, with a known and verified error.
So the project's three stages are not three separate exercises --- verification earns the
right to the measurement, measurement shows the object is interesting, and valuation shows it
is worth money.

\begin{saybox}[The sentence to end on]
``The model is wrong about the distribution of returns and still exactly right about the
accounting. That is not luck --- it is because the hedging identity is pathwise, and the only
thing it needs from the model is the quadratic variation, which is the one thing we verified
first.''
\end{saybox}

% =====================================================================
\section{``This looks simple'' --- the depth inventory}
\label{sec:depth}

A project can look simple because it is, or because it has been explained well. The difference
shows up the moment someone asks how a step works. This is the inventory of what is actually
non-trivial here, so that the answer exists.

\begin{table}[H]
\centering\small
\begin{tabular}{L{4.9cm}L{10.3cm}}
\toprule
\textbf{Step that looks routine} & \textbf{What is actually going on} \\
\midrule
``Volatility is the sum of squared returns'' &
This is a \emph{theorem}, not a definition: It\^o's lemma applied to $\log S$ plus the It\^o
isometry gives $[\log S]_T=\int\sigma_t^2\dd t$. A sum of squares estimates an SDE coefficient
only because Brownian paths have finite non-zero quadratic variation --- the same fact that
makes $\int\dd W$ impossible to define pathwise as a Stieltjes integral. \\
\addlinespace
``The error goes like $1/\sqrt n$'' &
We derive the \textbf{constant}, not just the rate:
$\operatorname{sd}=\sqrt{(2T/n)\int\sigma^4\dd t}$, via
$\operatorname{Var}((\Delta W)^2)=2\Delta^2$ from $(\Delta W)^2/\Delta\sim\chi^2_1$. Then we
match it to within $2.1\%$ over a $128\times$ range. Any estimator averaging independent terms
gives a $-1/2$ slope; almost none gives the right constant. \\
\addlinespace
``We delta-hedged'' &
The P\&L identity~\eqref{eq:pnl2} is derived from the self-financing condition plus It\^o plus
substitution of the Black--Scholes PDE, and the key step --- the $\dd W$ terms cancelling --- is
what makes the whole thing a \emph{variance} bet rather than a directional one. It is a
pathwise ODE solution, not an expectation. \\
\addlinespace
``We regressed realised on implied'' &
Mincer--Zarnowitz, with the joint null $a=0,b=1$, not merely $b\ne0$. And the choice of
\textbf{non-overlapping} windows costs us $95\%$ of the sample deliberately, because
overlapping windows would give badly understated standard errors. That is a real inferential
decision with a real price. \\
\addlinespace
``We checked for jumps'' &
Bipower variation, tripower quarticity, and the BNS ratio statistic with its
$(\pi/2)^2+\pi-5$ asymptotic constant --- a 2004--05 technique requiring an understanding of
why multiplying adjacent absolute returns is jump-robust while squaring is not. \\
\addlinespace
``We added transaction costs'' &
And discovered the drag grows as $n^{0.30}$, not $n^{1}$. The naive expectation is linear;
it is wrong because rebalancing more often shrinks each trade. Separating the deterministic
drag from the noisy mean is what makes that visible at all. \\
\addlinespace
``We used Black--Scholes'' &
With a verified implementation: PDE residual $3.4\times10^{-6}$ relative, Greeks against
central differences to $10^{-8}$, put--call parity to six decimals --- \emph{before} any result
was computed. \\
\bottomrule
\end{tabular}
\end{table}

\begin{saybox}[If the examiner says ``this seems straightforward'']
``The statement is simple --- volatility is a quadratic variation. Establishing that you can
actually measure it is not. I had to derive the estimator's asymptotic constant and check it to
$2\%$, prove the quadratic variation is $94\%$ diffusive rather than jumps, and show the
hedging identity holds pathwise on data with kurtosis 19. Each of those is a place the project
could have failed, and each one is a number in the report.''
\end{saybox}

\section{What every group member must be able to answer}
\label{sec:group}

A viva is marked on the group. The commonest failure is not a wrong answer from the presenter
but a blank one from a member who was not asked to understand a part they did not write. The
five items below are the minimum; nobody should go in without them.

\begin{keybox}[The five answers everyone needs]
\begin{enumerate}[leftmargin=*,topsep=1pt,itemsep=3pt]
\item \textbf{What is the project about, in one sentence?}\\
      ``It\^o's lemma says volatility is the quadratic variation of the log price, which makes
      it measurable from a single path. We verify that measurement, apply it to eleven years of
      NIFTY, and show the market prices volatility $2.5$ points too high.''
\item \textbf{Which numerical method was used, and which was not?}\\
      ``Exact log-Euler Monte Carlo --- exact because the GBM increment is closed-form in log
      coordinates --- closed-form Black--Scholes Greeks, a left-endpoint Riemann sum for the
      It\^o integral, and hand-written OLS. \textbf{No PDE solver, no tree, and no neural
      network or PINN anywhere.}''
\item \textbf{Where did the data come from?}\\
      ``Yahoo Finance via the \texttt{yfinance} client: NIFTY~50 (\texttt{\^{}NSEI}) and India
      VIX (\texttt{\^{}INDIAVIX}), 2{,}903 trading days from January 2015 to November 2026,
      cached locally so the numbers reproduce.''
\item \textbf{Why a 1973 model?}\\
      ``Black--Scholes is the null hypothesis, not the conclusion --- and we reject it. The
      identity we actually rely on, $[\log S]=\int\sigma^2\dd t$, is model-free, and the
      measurement tools are from 2001--2009.'' (Full answer: \S\ref{sec:1973}.)
\item \textbf{What is the headline number and what would break it?}\\
      ``Correlation $+0.949$ between the It\^o prediction and realised P\&L across 138 months.
      It would break if the price process had large jumps --- so we tested that, and $94\%$ of
      the quadratic variation is diffusive.''
\end{enumerate}
\end{keybox}

\begin{trapbox}[The specific way the last project went wrong]
In the Mathematical Modelling viva the group lost marks because members could not say which
method had been used where --- in particular whether a particular neural-network method had
been used at all. The fix is not more material; it is \textbf{one page that every member reads}.
That page is \S\ref{sec:methods}, and the single most important line in it is that
\emph{no neural method of any kind appears in this project}. If a member is unsure, the correct
answer is never ``I think so'' --- it is the sentence in item~2 above.
\end{trapbox}

\begin{keybox}[Suggested division for the viva]
Whoever presents should still be able to answer all five. But if the examiner directs questions
around the group, a natural split is: \textbf{theory and the derivations}
(\S\ref{sec:theory}); \textbf{data and validation} (\S\ref{sec:data}); \textbf{the verification
stage and its asymptotics}; \textbf{the measurement diagnostics and the jump test}; and
\textbf{the backtest, costs and limitations}. Each person should know their own area in depth
and the five answers above cold.
\end{keybox}

\section{What is genuinely weak}

Stating these yourself is worth more than having them extracted from you.

\begin{enumerate}[leftmargin=*,topsep=1pt,itemsep=3pt]
\item \textbf{The option premium is reconstructed, not quoted.} If he asks about the data,
      say this first. Both \emph{inputs} are real market data: NIFTY~50 closes, and an India
      VIX that the NSE itself builds from traded out-of-the-money NIFTY option prices by the
      variance-swap construction. But the option sold in the backtest is priced by
      \texttt{call\_price} --- Black--Scholes evaluated at $\si$ --- rather than taken from the
      quoted chain at that strike and expiry. The $+36.06$ points is therefore the premium the
      model says was available at the market's own implied volatility, not a fill anyone could
      demonstrate.

      \smallskip
      \textbf{What this does and does not contaminate.} Verification, measurement and the
      Mincer--Zarnowitz test use spot and VIX alone; nothing reconstructed touches them. The
      convergence slope, the Ljung--Box statistics, the leverage correlation, the jump test and
      the $b=0.664$ regression stand regardless. What inherits the model is the \emph{level} of
      the backtest P\&L --- not the variance gap $\si^2-\sigma_t^2$ driving it, which is
      measured from real data on both sides.

      \smallskip
      \textbf{Two things bound the error.} $\si$ is the model-free index rather than
      Black--Scholes inverted at a single strike, so the volatility input is the market's and
      not ours; and an at-the-money call is exactly where the skew correction to a
      Black--Scholes premium is smallest. \textbf{The honest next version is the NSE option
      chain with its bid--ask} --- volunteer that before he reaches for it.
\item \textbf{Costs are modelled, not measured.} The robustness section applies a proportional
      cost in basis points; it omits margin funding, market impact, and the bid--ask on the
      option itself. The sweep shows the premium survives to roughly $10$~bps round trip --- that
      is the claim, not that a desk would realise it net. \textbf{Do not claim an optimal
      rebalance frequency.} At zero cost the mean P\&L across frequencies is not statistically
      distinguishable (spread $22.5$ points against a typical standard error of $11.9$); what a
      higher frequency buys is lower dispersion, not a higher mean. Reading an optimum out of
      that table is reading noise, and that is the kind of overclaim that costs credibility on
      everything else in the room.
\item \textbf{Maturity mismatch.} India VIX is a 30-calendar-day index; our cycle is 21
      trading days ($\approx$30 calendar days, but not exactly). This biases the \emph{level}
      of $\si$ slightly. It does not affect the sign of the gap or the correlation result.
\item \textbf{Daily closes only.} $RV_n$ at daily frequency is coarse --- verification shows it is
      unbiased, but its one-month standard error is material. Intraday data would sharpen both
      the measurement and valuation considerably. Also, close-to-close sampling puts overnight jumps into $RV$ but
      not into the hedge, which cannot rebalance overnight.
\item \textbf{One market, one instrument.} At-the-money calls on a single index. The skew shows
      up in the leverage effect but is never actually priced.
\item \textbf{The $\sigma_r$ in the backtest is estimated, not known.} Unlike the verification
      stage, where we chose $\sigma(t)$, the realised volatility used in valuation is itself an
      $RV$ estimate with its own error. That error is quantified by the verification, which is
      why it comes first.
\end{enumerate}

% =====================================================================
\section{How to present this (15--20 minutes)}

The brief allots 15--20 minutes for the viva. Plan for about twelve minutes of material and
leave the rest for questions --- the questions are where the marks move.

\begin{table}[H]
\centering\small
\begin{tabular}{L{1.5cm}L{4.0cm}L{9.6cm}}
\toprule
\textbf{Time} & \textbf{Section} & \textbf{The line that carries it} \\
\midrule
0:00--1:30 & The problem &
``Black--Scholes needs five inputs. Four you can read off a screen. The fifth is $\sigma$, and
the course gives us exactly one tool for it: It\^o says $\sigma$ is a quadratic variation, so
it is measurable.'' \\
\addlinespace
1:30--4:00 & Verification &
``Before measuring anything real, check the ruler. We simulate with a $\sigma(t)$ we chose, so
we know the true answer. Theory predicts not just the rate but the constant --- and we match it
to $2\%$ over a $128\times$ range.'' \\
\addlinespace
4:00--6:00 & Measurement &
``Now point it at eleven years of NIFTY. Volatility moves between $5\%$ and $88\%$, clusters
enormously, and rises when the market falls. Constant-$\sigma$ GBM is wrong --- and we can say
exactly how.'' \\
\addlinespace
6:00--8:00 & Valuation &
``The market quotes its own volatility. Regress what follows on it: slope $0.66$, not $1$. The
market overreacts, by $2.5$ volatility points on average.'' \\
\addlinespace
8:00--10:00 & The payoff &
Derive \eqref{eq:pnl2} on the board --- four lines, pure course material --- then show the
scatter: $r=0.949$. \\
\addlinespace
10:00--12:00 & Robustness &
``Three things could break this, so we tested all three. Jumps: $94\%$ of the quadratic
variation is diffusive. Costs: the premium survives to $5$ basis points. Luck: every
sub-period is positive.'' \\
\addlinespace
12:00+ & Questions &
Lead with the two you want asked --- the 1973 question and the kurtosis question. Both have
prepared answers (\S\ref{sec:1973}, \S\ref{sec:defence}). \\
\bottomrule
\end{tabular}
\end{table}

\begin{keybox}[If you only remember six numbers]
$\mathbf{-0.4986}$ (the estimator converges as theory says) \ \textbullet\
$\mathbf{1752}$ vs $23.2$ (volatility clusters, overwhelmingly) \ \textbullet\
$\mathbf{0.664}$ (the market's forecast slope, not $1$) \ \textbullet\
$\mathbf{+0.949}$ (It\^o's prediction against reality) \ \textbullet\
$\mathbf{94\%}$ (share of quadratic variation that is diffusive, not jumps) \ \textbullet\
$\mathbf{+28.8}$ points (net premium at a realistic $2$~bps, $t=3.81$).
\end{keybox}

\begin{trapbox}[Three things to say before you are asked]
\textbf{1.} ``Black--Scholes here is the null hypothesis, and we reject it.'' \\
\textbf{2.} ``The returns are not normal --- excess kurtosis $19$ --- and the identity does not
need them to be, because it is pathwise.'' \\
\textbf{3.} ``The P\&L is net of modelled transaction costs, and it survives them.''
\smallskip
Volunteering all three converts your three biggest vulnerabilities into evidence that you
understand your own work. It also pre-empts roughly half the questions.
\end{trapbox}

\begin{qabox}[The questions most likely to come, and where the answer is]
\begin{itemize}[leftmargin=*,topsep=1pt,itemsep=1.5pt]
\item \emph{Why a model from 1973?} $\to$ \S\ref{sec:1973}, four-part answer.
\item \emph{Why not Heston / rough volatility?} $\to$ \S\ref{sec:1973}, final box.
\item \emph{Your returns aren't normal, so why does this work?} $\to$ \S\ref{sec:defence}.
\item \emph{Which numerical scheme, and why not Milstein?} $\to$ \S\ref{sec:methods}: in log
      coordinates the diffusion coefficient is state-independent, so Milstein reduces to Euler.
\item \emph{Did you use machine learning / PINNs?} $\to$ No. Nothing neural appears anywhere.
\item \emph{Isn't the sum of squared returns just a variance?} $\to$ \S\ref{sec:depth}, row 1:
      it estimates an SDE coefficient, and only because quadratic variation is finite and
      non-zero.
\item \emph{Where is this used in practice?} $\to$ \S\ref{sec:real}: variance swaps, India VIX
      futures, desk gamma P\&L attribution, margining.
\item \emph{Why only 138 observations?} $\to$ Non-overlapping windows, deliberately, so the
      standard errors are honest.
\item \emph{Is this a money-making strategy?} $\to$ No, and the report says so: it measures a
      risk premium, which is compensation for a real risk --- $25\%$ of months lose, the worst
      one by $419$ points.
\end{itemize}
\end{qabox}

\section{Reproducing everything}

\begin{lstlisting}[style=py,language=bash,numbers=none]
cd ~/Documents/SEM5/SC/01-Project
python sc_data.py      # fetch/cache + print sample statistics
python sc_model.py     # verify PDE residual, Greeks, put-call parity
python sc_vol.py       # verify, measure, value -> rv_convergence.csv, implied_vs_realised.csv
python sc_hedge.py     # convergence law + backtest -> convergence.csv, backtest_cycles.csv
python sc_robust.py    # jumps, costs, sub-periods  -> jump_by_year.csv, cost_sweep.csv
python sc_figures.py   # all six figures -> figures/*.pdf, *.png
latexmk -pdf SC_Report.tex
\end{lstlisting}

Requires \texttt{numpy}, \texttt{pandas}, \texttt{matplotlib} and \texttt{yfinance}; only the
first run of \texttt{sc\_data.py} needs a network connection, and only if \texttt{data/*.csv}
is absent.

\end{document}
